% ============================================================================
% CAPÍTULO 21 — Arquitetura\index{arquitetura} e escala
% Status: texto completo (rodada editorial 2)
% ============================================================================
\chapter{Arquitetura e Escala}
\label{cap:arquitetura}

\begin{caixaconceito}
Escalar não é ``fazer mais do mesmo'': é desenhar o sistema para crescer sem
desmoronar. Este capítulo cobre as peças que todo sistema full stack
produção encontra: \textbf{cache (Redis)}, \textbf{filas}, \textbf{bancos},
\textbf{microsserviços} e \textbf{padrões arquiteturais}.
\end{caixaconceito}

Até aqui, o SkillHub é uma aplicação única: um processo, um banco, tudo
síncrono. Isso é correto — e saudável. Mas chega o dia em que a aplicação
precisa fazer mais: enviar mil e-mails sem travar a requisição, servir a
mesma listagem para mil usuários sem bater no banco mil vezes, sobreviver a
picos sem desligar. Este capítulo é sobre esse dia: as decisões que você toma
\emph{antes} dele chegar.

Ao final deste capítulo, você será capaz de:
\begin{objetivos}
  \item Explicar os padrões de escala (vertical, horizontal, stateless\index{stateless});
  \item Usar Redis\index{Redis} para cache, sessões e rate limiting;
  \item Projetar filas\index{filas} assíncronas (processamento de tarefas em background);
  \item Comparar monólito\index{monólito}, modular e microsserviços\index{microsserviços} com critérios;
  \item Aplicar padrões: retry, circuit breaker, backoff, idempotência;
  \item Entregar o projeto do capítulo: a arquitetura de eventos do SkillHub.
\end{objetivos}

\section{Os três eixos de escala}

Quando um sistema fica lento, há três direções possíveis de crescimento — e
cada uma responde a uma pergunta diferente:

\begin{itemize}[leftmargin=*, itemsep=0.4em]
  \item \textbf{Vertical}: máquina maior (CPU/RAM). Simples, mas com teto e custo exponencial;
  \item \textbf{Horizontal}: mais instâncias atrás de um load balancer. Exige aplicação \emph{stateless} (estado fora do processo: banco/Redis);
  \item \textbf{Assíncrona}: tarefas pesadas saem da requisição e vão para filas — a requisição responde rápido e o processamento acontece em background.
\end{itemize}

\begin{caixaconceito}
\textbf{Stateless} é a chave: se cada instância guarda estado na memória
(sessões locais, cache local), adicionar instâncias quebra a aplicação. Estado
vai para serviços externos (PostgreSQL, Redis). É a diferença entre escalar e
``escalar até quebrar''.
\end{caixaconceito}

\begin{caixadica}
A escala vertical não é ``errada'' — é a primeira e mais barata. O erro é
tratá-la como a única: quando a máquina maior custa mais que duas máquinas
menores, o momento de ir horizontal chegou. E a escala assíncrona costuma ser
a mais subestimada: o e-mail que demora 2 segundos não precisa ser síncrono —
ele precisa \emph{sair da frente} da resposta HTTP.
\end{caixadica}

\section{Redis: o canivete suíço da infraestrutura}

O Redis é um banco em memória, chave-valor, absurdamente rápido (sub-
milissegundo). Ele não substitui o PostgreSQL — ele fica \emph{na frente}
dele para o que é lido muitas vezes e muda pouco:

\begin{lstlisting}[style=livro, language=TypeScript,
  caption={Cache com Redis (ioredis).}, label={list:arq-redis}]
import { Redis } from "ioredis";

const redis = new Redis(process.env.REDIS_URL!);

async function buscarTopServicos() {
  const CACHE_KEY = "skillhub:top-servicos";

  // 1. Tenta o cache
  const emCache = await redis.get(CACHE_KEY);
  if (emCache) return JSON.parse(emCache);

  // 2. Busca na fonte (lenta) e popula o cache com TTL
  const dados = await prisma.servico.findMany({ orderBy: { pedidos: "desc" }, take: 10 });
  await redis.set(CACHE_KEY, JSON.stringify(dados), "EX", 300); // 5 min
  return dados;
}
\end{lstlisting}

\begin{caixadica}
Padrão \textbf{cache-aside} (acima): leia do cache; em miss, busque na fonte e
popule com TTL. Cuidado com \emph{cache stampede} (muitas requisições
simultâneas em cache miss) — use lock ou \texttt{SET NX} no Redis.
\end{caixadica}

Além de cache, o Redis é a solução natural para três problemas de estado
compartilhado: \textbf{sessões} (o usuário logado em qualquer instância),
\textbf{rate limiting} (contador por IP/usuario com expiração — janela
deslizante) e \textbf{filas} (seção seguinte). São exatamente os casos em que
``estado na instância'' quebra com escala horizontal.

\section{Filas: processamento assíncrono}

O padrão é \textbf{produtor--fila--consumidor}: o request enfileira a tarefa e
responde \texttt{202 Accepted}; um worker processa em background.

\begin{itemize}[leftmargin=*, itemsep=0.4em]
  \item \textbf{BullMQ} (Redis): a fila padrão do ecossistema Node/TS — jobs, retries, agendamento, progresso;
  \item \textbf{Casos de uso}: e-mails, relatórios PDF, processamento de imagem, integrações lentas, notificações;
  \item \textbf{Idempotência}: processar a mesma mensagem duas vezes não pode duplicar efeito (chave única no job);
  \item \textbf{Dead letter}: jobs que falham N vezes vão para uma fila de inspeção humana.
\end{itemize}

\begin{lstlisting}[style=livro, language=TypeScript,
  caption={Produtor e consumidor com BullMQ.}, label={list:arq-fila}]
import { Queue, Worker } from "bullmq";

const filaEmail = new Queue("emails", { connection: redis });

// Produtor (dentro de uma rota)
await filaEmail.add("boas-vindas", { usuarioId: 42 }, {
  attempts: 3,
  backoff: { type: "exponential", delay: 2000 },
});

// Consumidor (processo separado)
new Worker("emails", async (job) => {
  const { usuarioId } = job.data;
  await enviarEmailBoasVindas(usuarioId); // lento, fora da requisição
}, { connection: redis });
\end{lstlisting}

\begin{caixaconceito}
\textbf{Por que a fila muda a arquitetura?} A requisição HTTP deixa de ser
responsável pelo resultado final da tarefa — ela só garante que a tarefa
\emph{foi aceita}. Isso desacopla tempo (a tarefa roda quando puder), volume
(a fila absorve picos) e falha (o retry cuida do erro sem travar o usuário).
O preço: o resultado deixa de ser imediato — e o produto precisa aceitar
``vai acontecer em instantes''.
\end{caixaconceito}

\section{Monólito vs microsserviços}

A decisão arquitetural mais discutida da década, resumida:

\begin{table}[h]
\centering
\small
\caption{Monólito, modular e microsserviços — critérios de decisão.}
\label{tab:arq-comparacao}
\begin{tabularx}{\textwidth}{@{}>{\raggedright\arraybackslash}p{2.8cm} >{\raggedright\arraybackslash}X >{\raggedright\arraybackslash}X >{\raggedright\arraybackslash}X@{}}
\toprule
\textbf{Critério} & \textbf{Monólito} & \textbf{Modular} & \textbf{Microsserviços} \\
\midrule
Deploys & 1 & 1 & N \\
Simplicidade & Alta & Alta & Baixa \\
Custo operacional & Baixo & Baixo & Alto \\
Escala independente & Não & Parcial & Sim \\
Time & Qualquer & Qualquer & Grande, por domínio \\
Quando escolher & Começo & A maioria & Caso específico \\
\bottomrule
\end{tabularx}
\end{table}

\begin{itemize}[leftmargin=*, itemsep=0.4em]
  \item \textbf{Monólito}: uma aplicação, um deploy. Simples, barato, fácil de começar — a escolha certa para 90\% dos produtos;
  \item \textbf{Modular (modulith)}: um deploy, mas módulos bem delimitados por domínio — o meio-termo recomendado;
  \item \textbf{Microsserviços}: serviços independentes por domínio, com seus próprios bancos e deploys. Justificável em times grandes e requisitos de escala independente — com custo alto de operação.
\end{itemize}

\begin{caixaconceito}
A regra prática: \textbf{comece monolítico, evolua com critério}. Extrair um
serviço (ex.: notificações) quando um domínio tiver requisitos de escala,
time e deploy próprios. Microsserviços por ``modernidade'' é a receita de
fracasso mais comum da indústria.
\end{caixaconceito}

\section{Padrões de resiliência}

Sistemas distribuídos falham — o design precisa assumir isso:

\begin{itemize}[leftmargin=*, itemsep=0.4em]
  \item \textbf{Retry com backoff}: exponencial + jitter para não sobrecarregar serviços lentos;
  \item \textbf{Circuit breaker}: após N falhas, ``abre o circuito'' e falha rápido por um período — protege o serviço degradado;
  \item \textbf{Timeout}: sempre defina timeout em chamadas externas (uma dependência lenta não pode travar seu servidor);
  \item \textbf{Idempotência}: chaves de idempotência em POST para evitar duplicação em retries;
  \item \textbf{Graceful degradation}: se a feature X cair, o resto continua funcionando.
\end{itemize}

\begin{caixaconceito}
\textbf{Circuit breaker} imita a engenharia elétrica: quando o serviço externo
está falhando, continuar chamando só piora (cada tentativa espera o timeout e
consome recursos). O breaker abre depois de N falhas — falha \emph{na hora},
sem chamar — e tenta de novo após um período (meio-aberto). O sistema degrada
com graça em vez de cair em cascata.
\end{caixaconceito}

% ============================================================================
\section{Projeto do capítulo: arquitetura de eventos do SkillHub}
\label{sec:projeto21}

\begin{caixaprojeto}
\textbf{Objetivo:} evoluir o SkillHub para processar notificações e
relatórios de forma assíncrona, com Redis e BullMQ.

\textbf{Critérios de aceite:}
\begin{itemize}[leftmargin=*, itemsep=0.2em]
  \item Redis integrado (Docker, capítulo~\ref{cap:docker}) para cache de listagem e sessões;
  \item Fila de e-mails com BullMQ: produtor na Server Action, worker separado, retry com backoff;
  \item Relatório mensal gerado em background (job agendado) e disponibilizado como arquivo;
  \item Rate limit do login usando Redis (janela deslizante);
  \item Cache-aside no ranking de serviços com invalidação sob demanda;
  \item Diagrama de arquitetura documentado (requisição síncrona vs fila assíncrona);
  \item Testes: produtor enfileira, worker processa, idempotência verificada.
\end{itemize}
\end{caixaprojeto}

\subsection{Passo a passo}

\begin{passos}
  \item \textbf{Redis no Compose}: adicione o serviço \texttt{redis} ao
  \texttt{docker-compose.yml} (capítulo~\ref{cap:docker}) com health check;
  \item \textbf{Cache-aside}: implemente o ranking de serviços com cache e
  TTL; invalide por tag ao criar um pedido;
  \item \textbf{Rate limit}: proteja a rota de login com contador Redis por
  e-mail/IP (janela deslizante) e mensagem genérica;
  \item \textbf{Fila de e-mails}: crie a fila BullMQ, o produtor na Server
  Action (e-mail de boas-vindas) e o worker em processo separado com retry
  exponencial;
  \item \textbf{Job agendado}: registre o worker de relatório mensal com
  \texttt{RepeatableJob} e publique o PDF gerado como arquivo;
  \item \textbf{Idempotência}: dê uma chave única a cada job (ex.:
  \texttt{email:boas-vindas:42}) para que reprocessamentos não dupliquem;
  \item \textbf{Documentação}: desenhe o diagrama síncrono vs assíncrono e
  escreva os testes de fila (produtor enfileira; worker processa; retry após
  falha).
\end{passos}

\section{Exercícios de fixação}

\begin{exercicios}
  \item Diferencie escala vertical, horizontal e assíncrona com um exemplo cada.
  \item Por que uma aplicação stateless é pré-requisito para escala horizontal?
  \item Desenhe o fluxo de uma tarefa em fila: quem produz, onde espera, quem consome.
  \item O que é idempotência e por que ela importa em filas e retries?
  \item Cite um caso em que microsserviços se justificam e um em que não.
\end{exercicios}

\section{Desafios}

\begin{desafios}
  \item \textbf{Circuit breaker}: implemente um circuit breaker simples para uma chamada externa (estado aberto/fechado/meio-aberto) e teste a recuperação.
  \item \textbf{Cache stampede}: simule 100 requisições simultâneas em cache miss e proteja com lock Redis; documente o ganho.
  \item \textbf{Fila de imagens}: processe redimensionamento de imagem em background (sharp) com fila e status de progresso.
\end{desafios}

\section{Exercícios de nível sênior}

\begin{seniores}
  \item \textbf{Event sourcing vs CRUD}: compare as duas abordagens para o domínio ``pedidos'' e descreva quando event sourcing se justifica.
  \item \textbf{Transações distribuídas}: explique por que transações ACID não cruzam microsserviços e como o padrão \emph{Saga} (coreografia vs orquestração) resolve.
  \item \textbf{Consistência eventual}: desenhe um caso em que cache + fila causam leitura defasada e documente a estratégia de reconciliação.
  \item \textbf{Backpressure}: projete o que acontece quando a fila enche (limite, drop com métricas, DLQ) e implemente a política.
\end{seniores}

\section{Armadilhas comuns}

\begin{itemize}[leftmargin=*, itemsep=0.4em]
  \item Cache sem TTL nem invalidação (dados velhos para sempre);
  \item Fila sem retry/backoff (uma falha de rede derruba o job para sempre);
  \item Microsserviços prematuros (custo operacional antes do problema real);
  \item Estado na instância (sessões locais) e depois ``não escala'';
  \item Sem timeouts em chamadas externas (cascata de lentidão);
  \item Job duplicado por falta de idempotência (o retry que deveria salvar, corrompe).
\end{itemize}

\section{Resumo do capítulo}

\begin{itemize}[leftmargin=*, itemsep=0.3em]
  \item Escala: vertical, horizontal (stateless) e assíncrona (filas);
  \item Redis: cache, sessões, rate limit, filas (BullMQ);
  \item Filas tiram trabalho pesado da requisição e ganham retry/backoff;
  \item Monólito primeiro; extraia serviços com critério;
  \item Resiliência: timeout, retry, circuit breaker, idempotência;
  \item Projeto entregue: SkillHub com arquitetura de eventos.
\end{itemize}

\section{Referências oficiais para aprofundar}

\begin{itemize}[leftmargin=*, itemsep=0.3em]
  \item Documentação do Redis: \url{https://redis.io/docs/latest/}
  \item BullMQ: \url{https://docs.bullmq.io}
  \item Microsoft — Padrões de nuvem (Cloud Design Patterns): \url{https://learn.microsoft.com/en-us/azure/architecture/patterns/}
  \item Martin Fowler — Microsserviços: \url{https://martinfowler.com/microservices/}
\end{itemize}
